% =====================================================================
%  REMAINING_ISSUES_BEFORE_AFTER.tex
%  Manuscript: Routing-Aware Federated Intrusion Detection with Robust
%  Aggregation and Median-Based Update Imputation
%  Source checked: main.tex (version of 2026-11-08 17:21, 18 pages)
%  Line numbers refer to that main.tex.
%  BEFORE = text as it stands now; AFTER = replacement LaTeX source.
%  [INSERT ...] = facts only the authors have; [VERIFY] = confirm first.
% =====================================================================
\documentclass[10pt,a4paper]{article}
\usepackage[T1]{fontenc}
\usepackage[utf8]{inputenc}
\usepackage[margin=2cm]{geometry}
\usepackage{xcolor}
\usepackage{fvextra}
\usepackage[most]{tcolorbox}
\usepackage{enumitem}
\usepackage[hidelinks]{hyperref}
\usepackage{lmodern}

\definecolor{bef}{HTML}{B3261E}
\definecolor{aft}{HTML}{1E7B34}
\definecolor{note}{HTML}{555555}

\fvset{breaklines=true,breakanywhere=true,fontsize=\small}

\newtcolorbox{before}[1][]{enhanced,breakable,colback=bef!4,colframe=bef!70,
  boxrule=0.6pt,arc=2pt,left=4pt,right=4pt,top=2pt,bottom=2pt,
  fonttitle=\bfseries\small,title={BEFORE #1}}
\newtcolorbox{after}[1][]{enhanced,breakable,colback=aft!4,colframe=aft!70,
  boxrule=0.6pt,arc=2pt,left=4pt,right=4pt,top=2pt,bottom=2pt,
  fonttitle=\bfseries\small,title={AFTER #1}}

\newcommand{\issue}[3]{\section{#1}{\raggedright\noindent\textcolor{note}{\small\textbf{Location:} #2\quad\textbf{Report item:} #3}\par}\smallskip}
\newcommand{\why}[1]{\noindent\textcolor{note}{\small #1}\par\smallskip}

\setlength{\parindent}{0pt}
\title{Remaining manuscript issues --- BEFORE / AFTER (LaTeX)}
\author{Checked against \texttt{main.tex}, 2026-11-08}
\date{}

\begin{document}
\maketitle
\tableofcontents
\bigskip

\noindent\textbf{Status.} Figures~1 and~3 and both captions are fixed (except the small wording point in Section~15). The items below are still open. Sections 1--15 can be applied now; Section~16 lists what needs data or author information.

% =====================================================================
\issue{Corresponding author and declarations}{lines 49--50; insert before line 2464 (\texttt{\textbackslash FloatBarrier})}{Comment 19}
\why{Elsevier/FGCS requires these statements; none is present.}

\begin{before}[lines 49--50]
\begin{Verbatim}
\author[1]{Nahideh Derakhshanfard}
\ead{n.derakhshan@iaut.ac.ir}
\end{Verbatim}
\end{before}

\begin{after}
\begin{Verbatim}
\author[1]{Nahideh Derakhshanfard\corref{cor1}}
\ead{n.derakhshan@iaut.ac.ir}
\cortext[cor1]{Corresponding author.}
\end{Verbatim}
\end{after}

\begin{before}[after the Data Availability section]
\begin{Verbatim}
(no declarations)
\end{Verbatim}
\end{before}

\begin{after}[insert before \texttt{\textbackslash FloatBarrier}, line 2464]
\begin{Verbatim}
\section*{CRediT authorship contribution statement}
\textbf{Nahideh Derakhshanfard:} [INSERT roles].
\textbf{Reza Azarnia:} [INSERT roles].
\textbf{Amir Mollanejad:} [INSERT roles].
\textbf{Ali Asghar Pour Haji Kazem:} [INSERT roles].
\textbf{Neda Dadashkhani:} [INSERT roles].
\textbf{Parisa Khoshvaght:} [INSERT roles].
\textbf{Mehdi Hosseinzadeh:} [INSERT roles].

\section*{Declaration of competing interest}
The authors declare that they have no known competing financial interests or
personal relationships that could have appeared to influence the work reported
in this paper. % [VERIFY]

\section*{Funding}
This research did not receive any specific grant from funding agencies in the
public, commercial, or not-for-profit sectors. % [VERIFY or INSERT agency + grant no.]

\section*{Ethics statement}
This study analyses publicly available network-traffic datasets (CICIoT2023 and
UNSW-NB15) and involves no human participants or animals.

\section*{Declaration of generative AI and AI-assisted technologies in the writing process}
During the preparation of this work the authors used [INSERT tool] in order to
[INSERT purpose]. After using this tool, the authors reviewed and edited the
content as needed and take full responsibility for the content of the publication.
\end{Verbatim}
\end{after}

% =====================================================================
\issue{Optimizer, loss, software and licence}{lines 1405--1406; line 2446}{Comment 18}
\why{\texttt{LocalOpt} is never specified; no licence is stated.}

\begin{before}[lines 1405--1406]
\begin{Verbatim}
The same base MLP architecture and common training settings were
used across configurations wherever applicable.
\end{Verbatim}
\end{before}

\begin{after}
\begin{Verbatim}
All configurations use the same MLP (hidden layers 128--96--64,
[INSERT activation] activations, dropout 0.2, one sigmoid output unit).
Local training minimizes binary cross-entropy with [INSERT optimizer]
(learning rate 0.0003, batch size 256, two local epochs); weights are
initialized with [INSERT scheme]. Experiments used Python [INSERT],
[INSERT framework and version], and the RDP accountant of
[INSERT library and version].
\end{Verbatim}
\end{after}

\begin{before}[line 2446, Code Availability]
\begin{Verbatim}
available in the \href{https://github.com/RezaAzarnia/routing-aware-fl-edgecloud-iot-ids}{GitHub repository}.
\end{Verbatim}
\end{before}

\begin{after}
\begin{Verbatim}
available in the \href{https://github.com/RezaAzarnia/routing-aware-fl-edgecloud-iot-ids}{GitHub repository}
under the [INSERT, e.g. MIT] licence.
\end{Verbatim}
\end{after}

% =====================================================================
\issue{Baselines that were never evaluated}{lines 1305--1309}{Comment 10}

\begin{before}[lines 1305--1309]
\begin{Verbatim}
In addition to standard federated-learning baselines, the experimental design
considers mechanisms motivated by modular and hierarchical resilient
federated learning approaches
\cite{almeida2025modular,chen2025hierarchical}.
\end{Verbatim}
\end{before}

\begin{after}
\begin{Verbatim}
The comparators are FedAvg under three selection/path settings
(TK/R, UR/R and UR/X), a non-private MAD-filtered reference and a
DP+FedAvg control (Table~\ref{tab:overall-results}). The modular and
hierarchical resilient frameworks of
\cite{almeida2025modular,chen2025hierarchical} were not
re-implemented, so no comparison with them is claimed.
\end{Verbatim}
\end{after}

% =====================================================================
\issue{Research gap and Table 1}{lines 231--241; Table 1, lines 508--649}{Comment 16}

\begin{before}[lines 231--241]
\begin{Verbatim}
Despite these advances, most existing studies address only part of the
overall problem. Some works mainly focus on intrusion detection and IDS
model quality, but they do not explicitly incorporate path stability
or communication failure into the learning process. Other studies
emphasize privacy preservation, secure aggregation, or Byzantine fault
tolerance while ignoring the impact of unstable paths, network
latency, and failed update delivery on FL decision-making.
Conversely, edge routing and scheduling approaches often aim to reduce
latency or energy consumption, but they rarely integrate robust
filtering for suspicious updates, central Gaussian privacy noise, and
failure recovery within a unified FL framework.
\end{Verbatim}
\end{before}

\begin{after}
\begin{Verbatim}
The studies closest to this work differ in which of four mechanisms they
evaluate together: routing-based path or client selection, handling of
missing updates, robust filtering of received updates, and a privacy
mechanism. Table~\ref{tab:related-work} records, for each study, whether
a mechanism is part of the method, reported as not evaluated, or absent.
None of the compared studies [VERIFY against Table 1] evaluates routing-
based selection together with missing-update imputation and MAD filtering
on an intrusion-detection benchmark; that combination is the subject of
this study.
\end{Verbatim}
\end{after}

\why{Table 1: replace the generic ``Main strength / Main limitation'' columns with mechanism columns and add the four systems named in the Introduction. Every \texttt{[V]} cell must be filled from the original paper.}

\begin{before}[Table 1 header and last row]
\begin{Verbatim}
\begin{tabularx}{\textwidth}{@{}l p{2.6cm} Y Y Y@{}}
Ref./year & Method & Main focus & Main strength & Main limitation \\
...
\textbf{Proposed}, 2026 & Routing-aware FL & IDS, routing, privacy, and recovery
  & Unified Edge--Cloud FL framework & Simulation-based; trusted or semi-trusted server \\
\end{Verbatim}
\end{before}

\begin{after}
\begin{Verbatim}
\begin{tabularx}{\textwidth}{@{}l p{2.2cm} Y Y Y Y Y@{}}
\toprule
Ref. & Method & Path/client selection & Missing updates & Robust filtering
  & Privacy mechanism & Evaluation data \\
\midrule
\cite{bhardwaj2024routing} & FL-based routing & Opportunistic routing & [V] & [V] & [V] & [V] \\
\cite{chen2025hierarchical} & Hierarchical FL & [V] & Fault tolerance [V] & [V] & [V] & [V] \\
\cite{geng2024smpc} & SMPC-based FL IDS & [V] & [V] & [V] & SMPC & [V] \\
\cite{lan2025ckks} & CKKS-based FL & [V] & [V] & [V] & CKKS encryption & [V] \\
\cite{soomro2026securedyn} & SecureDyn-FL & [V] & [V] & [V] & [V] & [V] \\
\cite{shah2026trisafe} & TriSAFE & [V] & [V] & [V] & Secure aggr.\ + DP & [V] \\
\cite{tang2026fedcst} & FedCST & [V] & [V] & [V] & [V] & [V] \\
\cite{alanazi2026edgeguard} & EdgeGuard-AI & Load-aware scheduling & [V] & [V] & [V] & [V] \\
% ... remaining references, same columns ...
\textbf{This study} & Routing-aware FL & Top-$k$ by delay, stability, security
  & Median surrogate (no information recovered) & MAD
  & Server-side Gaussian; diagnostic RDP, not certified
  & CICIoT2023, UNSW-NB15 (simulated network) \\
\bottomrule
\end{tabularx}
\end{Verbatim}
\end{after}

% =====================================================================
\issue{``Recovery'' wording for the surrogate (text no longer matches Figure 1)}{20 places, listed below}{Comment 5}
\why{Figure 1 now says ``Surrogate Imputation''. ``Recovery'' stays only where it means F1 recovery (Table 6, recovery rounds). Also rename every \texttt{\textbackslash eqref\{eq:recovered-update\}}.}

\begin{before}[one line per location]
\begin{Verbatim}
L241   failure recovery within a unified FL framework.
L643   & IDS, routing, privacy, and recovery
L666   recovery, filtering, privacy-aware aggregation, and global-model
L675   updates and applies robust filtering and surrogate recovery, with
L832   received updates using failure recovery, robust filtering, and
L834   \subsection{Failure Recovery and Robust Update Filtering}
L846   missing-update surrogates are reconstructed before MAD filtering.
L849   reconstructed from the retained real updates.
L853-854 a surrogate update is generated using coordinate-wise median reconstruction:
L868   \label{eq:recovered-update}      (and \eqref at L884)
L891   After recovery, the server constructs the available update set
L893   and recovered surrogate updates. The robust center and the deviation
L967   Following robust filtering and any configured surrogate recovery,
L1239  Reconstruct missing updates using the coordinate-wise median
L1315  behavior, robust filtering, surrogate recovery, weighting, and
L1716-1718 The update reconstruction ratio measures the proportion of failed or
       missing model updates successfully reconstructed by the recovery mechanism:
L1723  \frac{N_{\mathrm{recovered}}}{N_{\mathrm{failed}}+\xi}
L1822  Recovered-update coefficient
L1849  Failure recovery & --- & Coordinate-wise median \\
L2185  changes for the remaining controls. Removing surrogate recovery,
L2188  surrogate reconstruction did not produce clearly distinguishable
\end{Verbatim}
\end{before}

\begin{after}
\begin{Verbatim}
L241   missing-update imputation within the same FL framework.
L643   & IDS, routing, privacy, and update imputation
L666   surrogate imputation, filtering, aggregation, and global-model
L675   updates and applies robust filtering and surrogate imputation, with
L832   received updates using surrogate imputation, robust filtering, and
L834   \subsection{Surrogate Imputation and Robust Update Filtering}
L846   missing-update surrogates are computed before MAD filtering.
L849   computed from the retained real updates.
L853-854 a surrogate update is generated using coordinate-wise median imputation:
L868   \label{eq:surrogate-update}      (and \eqref{eq:surrogate-update} at L884)
L891   After imputation, the server constructs the available update set
L893   and surrogate updates. The robust center and the deviation
L967   Following robust filtering and any configured surrogate imputation,
L1239  Impute missing updates with the coordinate-wise median
L1315  behavior, robust filtering, surrogate imputation, weighting, and
L1716-1718 The surrogate availability ratio measures the proportion of failed or
       missing model updates for which a surrogate was generated:
L1723  \frac{N_{\mathrm{imputed}}}{N_{\mathrm{failed}}+\xi}   (also L1728)
L1822  Surrogate raw-weight factor
L1849  Missing-update handling & --- & Coordinate-wise median surrogate \\
L2185  changes for the remaining controls. Removing surrogate imputation,
L2188  surrogate imputation did not produce clearly distinguishable
\end{Verbatim}
\end{after}

% =====================================================================
\issue{Metrics defined but never reported}{lines 1608--1631 and 1678--1695}{Comment 13}
\why{Eq.~(24) Recovery Time and Eq.~(27) Communication cost have no reported values. Delete Eq.~(24) and its paragraph (lines 1613--1631), and pin Eq.~(27) to the scalability results.}

\begin{before}[lines 1608--1609]
\begin{Verbatim}
Resilience was evaluated using recovery time, recovery rounds, and
recovery success.
\end{Verbatim}
\end{before}

\begin{after}
\begin{Verbatim}
Resilience was evaluated using the F1 recovery fraction, recovery
rounds, and the surrogate availability ratio.
\end{Verbatim}
\end{after}

\begin{before}[lines 1678--1695]
\begin{Verbatim}
Communication cost measures the total volume of model-update data
transmitted by participating clients during federated training:
  (Eq. communication-cost)
A lower communication cost is generally preferable. However, it
should be interpreted together with end-to-end delay, PDR, path
stability, transfer efficiency, and intrusion-detection performance.
\end{Verbatim}
\end{before}

\begin{after}
\begin{Verbatim}
Communication cost is the simulated serialized uplink traffic of the
selected clients, summed over rounds:
  (Eq. communication-cost, unchanged)
Each update holds [INSERT P] parameters stored as [INSERT float32]
([INSERT] MB per update); downlink broadcast, metadata,
retransmissions and imputed surrogates are not counted. It is
reported for the scalability study (Figure~\ref{fig:scalability-new}).
\end{Verbatim}
\end{after}

% =====================================================================
\issue{Zero-weight surrogate sentence contradicts the code and Figure 1}{lines 1052--1058}{new (found in code audit)}
\why{\texttt{capped\_normalized\_weights} gives every retained update, including a zero-raw-weight surrogate, exactly $1/|\mathcal K_t|$ whenever the raw-weight sum is positive. The sentence is only true for the abstention case.}

\begin{before}[lines 1052--1058]
\begin{Verbatim}
When all retained raw weights sum to zero, the SUPP\_V2 corrective
implementation takes a \emph{no-valid-contribution} abstention:
the global model and server momentum remain unchanged, no Gaussian
noise is released, and no zero-weight surrogate gains positive
weight. Release-based and attempted-round accounting diagnostics
are recorded separately. This branch does not change any
positive-weight aggregation round.
\end{Verbatim}
\end{before}

\begin{after}
\begin{Verbatim}
When all retained raw weights sum to zero, the SUPP\_V2 corrective
implementation takes a \emph{no-valid-contribution} abstention:
the global model and server momentum remain unchanged and no Gaussian
noise is released. Release-based and attempted-round accounting
diagnostics are recorded separately. This branch does not change any
positive-weight aggregation round; in such rounds the uniform cap
assigns $1/|\mathcal K_t|$ to every retained update, including a
surrogate whose raw weight is zero.
\end{Verbatim}
\end{after}

% =====================================================================
\issue{CapNormalize has no formula}{lines 1033--1035}{Comment 6}

\begin{before}[lines 1033--1035]
\begin{Verbatim}
For a round with a positive raw-weight sum, the central-Gaussian
weights satisfy \(\sum_{i\in\mathcal K_t}a_i^{(t)}=1\) and
\(0\leq a_i^{(t)}\leq\overline a_{\max}^{(t)}\).
\end{Verbatim}
\end{before}

\begin{after}
\begin{Verbatim}
\(\operatorname{CapNormalize}\) assigns
\(a_i^{(t)}=\min\bigl(\overline a_{\max}^{(t)},\,\lambda_t b_i^{(t)}\bigr)\),
iteratively fixing capped entries and distributing the remaining mass
among the others, with \(\lambda_t>0\) chosen so that the weights sum
to one. For a round with a positive raw-weight sum, the central-Gaussian
weights therefore satisfy \(\sum_{i\in\mathcal K_t}a_i^{(t)}=1\) and
\(0\leq a_i^{(t)}\leq\overline a_{\max}^{(t)}\).
\end{Verbatim}
\end{after}

% =====================================================================
\issue{Adjacency and accountant details}{after line 1109; Table 3, line 1801}{Comment 3}

\begin{before}[line 1108]
\begin{Verbatim}
The implementation tracks the cumulative privacy budget using an
RDP accountant.
\end{Verbatim}
\end{before}

\begin{after}
\begin{Verbatim}
The implementation tracks the cumulative privacy budget using an
RDP accountant ([INSERT library and version], R\'enyi orders
[INSERT orders]). Adjacency is defined at the [INSERT client / record]
level, i.e.\ adding or removing one [INSERT client's update / record].
\end{Verbatim}
\end{after}

\begin{before}[line 1801, Table 3]
\begin{Verbatim}
	Target privacy budgets
\end{Verbatim}
\end{before}

\begin{after}
\begin{Verbatim}
	Diagnostic RDP targets
\end{Verbatim}
\end{after}

% =====================================================================
\issue{Network model: link dynamics and failure combination}{lines 1383--1384}{Comment 4}

\begin{before}[lines 1382--1384]
\begin{Verbatim}
Link attributes
vary stochastically by round, and client--edge outages occur with
probability 0.01.
\end{Verbatim}
\end{before}

\begin{after}
\begin{Verbatim}
Link attributes vary by round as [INSERT process, e.g. a bounded
Gaussian perturbation with \(\sigma=0.025\)] [VERIFY], and client--edge
outages occur independently with probability 0.01. Injected client
failures (\(p_f\)) are drawn independently of link outages; an update
is missing if either event occurs. [VERIFY]
\end{Verbatim}
\end{after}

% =====================================================================
\issue{Data provenance}{line 1374}{Comment 8}

\begin{before}[lines 1373--1374]
\begin{Verbatim}
The split and
preparation fingerprints are documented in the locked manifests.
\end{Verbatim}
\end{before}

\begin{after}
\begin{Verbatim}
The split and
preparation fingerprints are documented in the locked manifests.
Attack/benign counts are [INSERT] (CICIoT2023 training),
[INSERT] (CICIoT2023 validation), [INSERT] (UNSW-NB15 training) and
[INSERT] (UNSW-NB15 validation). Scaling and categorical encoders are
fitted on training records only. Client partitions draw, for each class,
client proportions from a Dirichlet distribution with
\(\alpha_{\mathrm{Dir}}=0.5\) [VERIFY], with a minimum of 100 records per
client.
\end{Verbatim}
\end{after}

% =====================================================================
\issue{Abstract: multi-factor comparison stated without qualification}{lines 102--105}{Comment 17}

\begin{before}[lines 102--105]
\begin{Verbatim}
Relative to FedAvg with uniform selection and
random paths, its mean delivery delay decreases from 87.3 to
66.0~ms, while first-transmission PDR increases from 0.694 to
0.851.
\end{Verbatim}
\end{before}

\begin{after}
\begin{Verbatim}
Relative to FedAvg with uniform selection and
random paths, a comparison that changes selection, routing and
aggregation together, its mean delivery delay decreases from 87.3 to
66.0~ms, while first-transmission PDR increases from 0.694 to
0.851.
\end{Verbatim}
\end{after}

% =====================================================================
\issue{Contributions and keywords overclaim}{lines 270--276; lines 119--128}{Comment 20}

\begin{before}[lines 270--276]
\begin{Verbatim}
\item An integrated FL framework is proposed to jointly address
intrusion detection, routing-aware communication, privacy-aware
aggregation, and fault tolerance in Edge--Cloud IoT systems.

\item A routing-aware client and path selection scheme is designed
using latency, path stability, and route-security indicators to
reduce end-to-end delay and improve communication reliability.
\end{Verbatim}
\end{before}

\begin{after}
\begin{Verbatim}
\item An FL pipeline for intrusion detection is implemented in which
routing-aware selection, surrogate imputation, MAD filtering and
optional server-side Gaussian perturbation can be switched on and off,
so that each can be evaluated in controlled comparisons.

\item Clients and paths are ranked by simulated latency, stability and
security indicators; the effect of this ranking on simulated delivery
delay and packet delivery is measured separately from its effect on
detection.
\end{Verbatim}
\end{after}

\begin{before}[lines 119--128]
\begin{Verbatim}
\begin{keyword}
Federated learning \sep Internet of Things \sep Edge--Cloud architecture
\sep Intrusion detection \sep Routing-aware aggregation \sep Fault tolerance
\sep Privacy-aware aggregation \sep Network resilience
\end{keyword}
\end{Verbatim}
\end{before}

\begin{after}
\begin{Verbatim}
\begin{keyword}
Federated learning \sep IoT intrusion detection \sep Routing-aware client selection
\sep Robust aggregation \sep Surrogate update imputation
\sep Gaussian perturbation \sep Communication failures
\end{keyword}
\end{Verbatim}
\end{after}

% =====================================================================
\issue{Generic sentences in the Introduction and section openings}{lines 134--146, 149--151, 171--175, 201--203, 307--315, 654--656, 1298, 1338--1341}{Comment 20 / Section D}

\begin{before}[lines 134--137]
\begin{Verbatim}
The rapid expansion of Internet of Things (IoT) systems in smart
cities, digital healthcare, intelligent manufacturing, transportation,
and critical infrastructure has generated massive volumes of
distributed, heterogeneous, and highly sensitive data.
\end{Verbatim}
\end{before}
\begin{after}
\begin{Verbatim}
IoT intrusion-detection models are increasingly trained across many
devices whose traffic data cannot be pooled centrally.
\end{Verbatim}
\end{after}

\begin{before}[lines 149--151]
\begin{Verbatim}
learning paradigm that enables IDS models to be trained without
transferring raw data, thereby reducing the risk of information
exposure.
\end{Verbatim}
\end{before}
\begin{after}
\begin{Verbatim}
learning paradigm that enables IDS models to be trained without
transferring the training records, although transmitted model updates
may still expose information to the receiving server.
\end{Verbatim}
\end{after}

\begin{before}[lines 171--175]
\begin{Verbatim}
consider privacy and Byzantine resilience indicate that FL security
should go beyond decentralized model training and include client
behavior analysis, reliability assessment, and robust aggregation
\cite{shu2026trustworthy}.
\end{Verbatim}
\end{before}
\begin{after}
\begin{Verbatim}
consider privacy and Byzantine resilience treat update confidentiality,
poisoning resistance and delivery reliability as separate requirements,
each needing its own assumptions and evaluation
\cite{shu2026trustworthy}.
\end{Verbatim}
\end{after}

\begin{before}[lines 201--203]
\begin{Verbatim}
Although these approaches can provide strong
confidentiality and secure aggregation, they often introduce
considerable computational or communication overhead.
\end{Verbatim}
\end{before}
\begin{after}
\begin{Verbatim}
These approaches add cryptographic or protocol steps whose computation
and communication costs are reported in [VERIFY: cite the figures from
the original papers].
\end{Verbatim}
\end{after}

\begin{before}[lines 307--315]
\begin{Verbatim}
This section reviews federated-learning-based studies related to
IoT security, fault tolerance, privacy preservation, and routing
optimization in Edge--Cloud IoT systems. The reviewed studies are
organized into four main categories: ... These research directions are
reviewed to identify the remaining gaps addressed by the proposed
framework.
\end{Verbatim}
\end{before}
\begin{after}
\begin{Verbatim}
The studies below are grouped by the mechanism they contribute: IDS
model quality, failure handling, privacy protection, and path or client
selection. Table~\ref{tab:related-work} compares them on those four
mechanisms.
\end{Verbatim}
\end{after}

\begin{before}[lines 654--656]
\begin{Verbatim}
This study proposes a federated learning framework with
routing-aware, privacy-aware, and fault-tolerant mechanisms for
intrusion detection in Edge--Cloud IoT systems.
\end{Verbatim}
\end{before}
\begin{after}
\begin{Verbatim}
This study evaluates a federated learning pipeline that combines
routing-aware selection, surrogate imputation, robust filtering and
optional Gaussian perturbation for intrusion detection in Edge--Cloud
IoT systems.
\end{Verbatim}
\end{after}

\begin{before}[line 1298]
\begin{Verbatim}
\section{Simulation}
\end{Verbatim}
\end{before}
\begin{after}
\begin{Verbatim}
\section{Experimental Evaluation}
\end{Verbatim}
\end{after}

\begin{before}[lines 1338--1341]
\begin{Verbatim}
This experimental design enables the proposed framework to be
evaluated not only in terms of intrusion detection performance, but
also with respect to routing stability, communication delay, privacy
awareness, and fault tolerance.
\end{Verbatim}
\end{before}
\begin{after}
\begin{Verbatim}
% (delete: repeats the opening of Section 4)
\end{Verbatim}
\end{after}

% =====================================================================
\issue{Figure 3 caption wording}{lines 2026--2031}{Comment 12 (follow-up)}
\why{``Gaussian and adaptive attacks'' reads as if a ``Gaussian attack'' existed; the point is that the Gaussian-perturbed configuration was not tested. The methods are now also distinguished by marker shape.}

\begin{before}[lines 2026--2031]
\begin{Verbatim}
\caption{Non-adaptive attack comparison using paired seeds 12--16
	on UNSW-NB15 and CICIoT2023. Both arms use top-$k$ selection and
	routed paths: orange, unfiltered FedAvg; blue, non-private MAD with
	median imputation. Top/bottom panels show selected/last-round ROC--AUC;
	dots indicate seeds and bars means. SF: sign flip; BSF: norm-preserving
	sign flip; LF: label flip. Gaussian and adaptive attacks are not tested.}
\end{Verbatim}
\end{before}

\begin{after}
\begin{Verbatim}
\caption{Non-adaptive attack comparison using paired seeds 12--16
	on UNSW-NB15 and CICIoT2023. Both arms use top-$k$ selection and
	routed paths: orange circles, unfiltered FedAvg; blue squares,
	non-private MAD with median imputation. Top/bottom panels show
	selected/last-round ROC--AUC; markers indicate seeds and bars means.
	SF: sign flip; BSF: norm-preserving sign flip; LF: label flip.
	The Gaussian-perturbed configuration and adaptive attacks were not
	evaluated.}
\end{Verbatim}
\end{after}

% =====================================================================
\section{Items that need data or author information}
These cannot be written from the manuscript alone:
\begin{enumerate}[leftmargin=*,itemsep=2pt]
\item \textbf{Reserved test sets were never reported.} Run the frozen-threshold final evaluation of the main Table~4 rows, or state explicitly why it is withheld.
\item Confusion matrices, per-seed thresholds, and recall at 1\% and 5\% validation FPR (Comment 9).
\item An explanation for $\mathrm{AUC}_{\mathrm{last}}=0.5000\,(0.0000)$ of FedAvg (UR/X) on CICIoT2023 (Table~4).
\item Line 2264: ``server learning rate 5.0'' versus 1.0 in the review report --- check against the logs.
\item Longest participation gap per policy (Comment 7); the column \texttt{max\_selection\_gap} already exists in the supplementary per-seed CSV.
\item Attack/benign counts per split; parameter count and bytes per update.
\item CRediT roles, funding, competing interests, AI-use statement, optimizer, activation, initialization, software versions, licence, adjacency, accountant library and orders.
\item Hyperparameter-selection rationale ($k=2.5$, $C$, $\chi_{\mathrm{rec}}$) and tuning budget (Comment 15).
\item All \texttt{[V]} cells of the new Table~1, taken from the original papers.
\item A new Zenodo version that includes SUPP\_V1/V2, with the DOI updated in Code/Data Availability; submit Supplementary Figure~S2 and Table~S6.
\item DOIs for references that lack them (e.g.\ [1]--[3], [5]--[9]).
\end{enumerate}

\end{document}
